\documentclass[10pt,a4paper]{amsart}
\usepackage[margin=19mm]{geometry}
\usepackage{amsmath,amssymb,mathtools}
\usepackage[scr=rsfs,cal=euler]{mathalfa}
\usepackage{xfrac}
\usepackage[unicode,hidelinks,bookmarks=false]{hyperref}
\newcommand{\ie}{i.e.\ }
\newcommand{\cf}{cf.\ }
\newcommand{\resp}{respectively\ }
\newcommand{\Subsection}{\subsection}
\providecommand{\nobreakdash}{\nobreak\hbox{-}}
\setlength{\emergencystretch}{2em}
\multlinegap0pt
\usepackage{bm}
\usepackage{bbm}
\usepackage{dsfont}
\usepackage{xspace}
\usepackage{cancel}
\usepackage{mathtools}
\usepackage{amsthm}
\makeatletter
\@ifundefined{theorem}{\newtheorem{theorem}{Theorem}}{}
\@ifundefined{example}{\newtheorem{example}[theorem]{Example}}{}
\makeatother
\title[Hypersurfaces of $\mathbb{H}^2\times\mathbb{H}^2$ with const]{Hypersurfaces of $\mathbb{H}^2\times\mathbb{H}^2$ with constant sectional curvature}
\author{Haizhong Li, Luc Vrancken, Xianfeng Wang, Zeke Yao}
\date{Source: arXiv:2606.24686v1 (CC BY 4.0)}
\begin{document}
\maketitle

\noindent\emph{Source and licence.} Hypersurfaces of $\mathbb{H}^2\times\mathbb{H}^2$ with constant sectional curvature. Version arXiv:2606.24686v1 (CC BY 4.0), distributed under the Creative Commons Attribution 4.0 International licence, as recorded in the arXiv licence metadata for this exact version and shown on the abstract page https://arxiv.org/abs/2606.24686v1. Full rights notice: licenses/arxiv-2606.24686-CC-BY-4.0.txt.\par\smallskip

\noindent\textit{Editorial scope.} Counted unit: the Theorem whose source label is \texttt{thm:4.1abc} in the article (source0.tex lines 1940--1959, source label \texttt{thm:4.1abc}), delivered together with the article's own complete original proof of it (source0.tex lines 1961--2221, one \texttt{proof} environment of 261 lines). No mathematics is written, repaired or re-derived: statement and proof are the article's own text line for line. Required context retained uncounted: eqn:2.2 (lines 461--467); eqn:2.5 (lines 628--631); eqn:3.1ss1 (lines 776--781); eqn:3.3ss1 (lines 823--830); eqn:3.4ss1 (lines 833--841); exam:4.1ss1ss11ab1 (lines 905--930); eqn:jr (lines 920--927); eqn:LL (lines 1803--1817). Every definition, display and label of those lines is retained unchanged. Omitted from this package and not redistributed: 1--460, 468--627, 632--775, 782--822, 831--832, 842--904, 928--1802, 1818--1939, 1960--1960, 2222--4280. Nothing inside a retained excerpt is omitted or altered: every retained body is delivered exactly as the article's lines are written, so no omission receipt of any kind accompanies this package. Citations: the retained excerpts carry no citation command at all, so no bibliography entry of the article is transcribed and no reference is redistributed. Reference adaptation: none was needed and none was made; every reference inside the delivered unit resolves inside it, so no unresolved reference or citation is produced. Printed slips kept literal: no printed word, symbol, hypothesis or number of the counted statement or of its proof was altered. Layout and class: the article's own class is \texttt{amsart} with packages including babel, amsmath, amssymb, amsthm, hyperref, geometry, xcolor; the wrapper is the lane's pinned \texttt{amsart} editorial wrapper at 10pt on A4, which restates the theorem-like environments the retained text needs and adds the article's own macro definitions that the retained text needs (none), with extra wrapper packages (bm, bbm, dsfont, xspace, cancel, mathtools). The delivered numbering is the unit's own. Assets: no retained span contains a figure environment, a graphics-inclusion command, a PDF-page inclusion, a table or a TikZ picture; no image file is copied into this package, the package declares no asset dependency and no image file is redistributed with it. Licence: version arXiv:2606.24686v1 of this article is distributed under the Creative Commons Attribution 4.0 International licence, as recorded in the arXiv licence metadata for this exact version and shown on the abstract page https://arxiv.org/abs/2606.24686v1. A full rights notice is delivered at licenses/arxiv-2606.24686-CC-BY-4.0.txt. Characters: every retained excerpt is pure ASCII, so the delivered engine, XeLaTeX with its default Latin Modern text font, reports no missing character.
\begin{equation}\label{eqn:2.2}
\begin{aligned}
\operatorname{R}(U, Y)Z=&-\tfrac{1}{2}\big[g(Y,Z)U-g(U,Z)Y+g(TY,Z)TU
-g(TU,Z)TY\big]\\
&+g(AY,Z)AU-g(AU,Z)AY,
\end{aligned}
\end{equation}

\begin{equation}\label{eqn:2.5}
E_1=\frac{J_1N+J_2 N}{\sqrt{2(1+C)}},\
\ E_2=\frac{J_1N-J_2 N}{\sqrt{2(1-C)}},\ \ E_3=\frac{X}{\sqrt{1-C^{2}}}.
\end{equation}

\begin{equation}\label{eqn:3.1ss1}
\begin{aligned}
&p(t,r)=\cosh(f_1(t))\gamma_1(r)+\sinh(f_1(t))N_1(r),\\
&q(t,s)=\cosh(f_2(t))\gamma_2(s)+\sinh(f_2(t))N_2(s),
\end{aligned}	
\end{equation}

\begin{equation}\label{eqn:3.3ss1}
\begin{aligned}
N&=\frac{1}{\sqrt{(f_1'(t))^2+(f_2'(t))^2}}(f_2'(t)\sinh(f_1(t))\gamma_1(r)
+f_2'(t)\cosh(f_1(t))N_1(r),\\
&\qquad \qquad -f_1'(t)\sinh(f_2(t))\gamma_2(s)
-f_1'(t)\cosh(f_2(t))N_2(s)).
\end{aligned}
\end{equation}

\begin{equation}\label{eqn:3.4ss1}
\begin{aligned}
&AE_1=-\frac{f_2'(t)}{\sqrt{(f_1'(t))^2+(f_2'(t))^2}}\frac{\sinh(f_1(t))-k_1(r)\cosh(f_1(t))}
{\cosh(f_1(t))-k_1(r)\sinh(f_1(t))}E_1,\\
&AE_2=\frac{f_1'(t)}{\sqrt{(f_1'(t))^2+(f_2'(t))^2}}\frac{\sinh(f_2(t))-k_2(s)\cosh(f_2(t))}
{\cosh(f_2(t))-k_2(s)\sinh(f_2(t))}E_2,\\
&AE_3=-\frac{f_1'(t)f_2''(t)-f_1''(t)f_2'(t)}{((f_1'(t))^2+(f_2'(t))^2)^{\frac{3}{2}}}E_3.
\end{aligned}
\end{equation}

\begin{example}\label{exam:4.1ss1ss11ab1}
In \eqref{eqn:3.1ss1}, for some constant $0\neq a\in \mathbb{R}$, we take
$$
\begin{aligned}
&f_1(t)=\ln\Big(\sqrt{2}\cosh(\tfrac{t}{\sqrt{2}})+\sqrt{\cosh(\sqrt{2}t)}\Big),\ \
f_2(t)={\rm arcsinh}\Big(-\sqrt{2}\sinh(\tfrac{t}{\sqrt{2}})\Big)-a,\ \ t<0,\\
&\gamma_1(r)=(\cosh(r),\sinh(r),0), \ \
\gamma_2(s)=(\cosh(a),\sinh(a)\cos(s),\sinh(a)\sin(s)),\\
&N_1(r)=(0,0,1),\ \ N_2(s)=(\sinh(a),\cosh(a)\cos(s),\cosh(a)\sin(s)).
\end{aligned}	
$$
Then we have the immersion:
$\Phi_1:\Omega\subset\mathbb{R}^3\rightarrow\mathbb{H}^2\times\mathbb{H}^2$:
$(t,r,s)\rightarrow(p(t,r),q(t,s))$, where %$t<0$ and
%
\begin{equation}\label{eqn:jr}
\begin{aligned}
%&p(t,r)=\sqrt{\cosh(\sqrt{2}t)+1}(\cosh(r),\sinh(r),0)+\sqrt{\cosh(\sqrt{2}t)}(0,0,1),\\
%&q(t,s)=\sqrt{\cosh(\sqrt{2}t)}(1,0,0)+\sqrt{\cosh(\sqrt{2}t)-1}(0,\cos(s),\sin(s)),
&p(t,r)=\Big(\sqrt{\cosh(\sqrt{2}t)+1}\cosh(r),\sqrt{\cosh(\sqrt{2}t)+1}\sinh(r),\sqrt{\cosh(\sqrt{2}t)}\Big),\\
&q(t,s)=\Big(\sqrt{\cosh(\sqrt{2}t)},\sqrt{\cosh(\sqrt{2}t)-1}\cos(s),\sqrt{\cosh(\sqrt{2}t)-1}\sin(s)\Big).
\end{aligned}	
\end{equation}
By \eqref{eqn:2.2}, \eqref{eqn:3.3ss1} and \eqref{eqn:3.4ss1}, one can be checked that such hypersurface constructed by \eqref{eqn:jr} has constant sectional curvature $-\frac{1}{2}$, and
$C={\rm sech}(\sqrt{2}t)$.
\end{example}

\begin{equation}\label{eqn:LL}
\left\{
\begin{aligned}
&E_1C=-2b_3\sqrt{1-C^2},\ E_2C=-2b_5\sqrt{1-C^2},\ E_3C=-2b_6\sqrt{1-C^2},\\ 		
&\nabla_{E_1}E_1=b_1\sqrt{\tfrac{1-C}{1+C}}E_3,\
\nabla_{E_1}E_2=-b_2\sqrt{\tfrac{1+C}{1-C}}E_3,\ \nabla_{E_1}E_3=-b_1\sqrt{\tfrac{1-C}{1+C}}E_1
+b_2\sqrt{\tfrac{1+C}{1-C}}E_2,\\
&\nabla_{E_2} E_1=b_2\sqrt{\tfrac{1-C}{1+C}}E_3,\
\nabla_{E_2} E_2=-b_4\sqrt{\tfrac{1+C}{1-C}}E_3,\ \nabla_{E_2}E_3=-b_2\sqrt{\tfrac{1-C}{1+C}}E_1
+b_4\sqrt{\tfrac{1+C}{1-C}}E_2,\\
&\nabla_{E_3}E_1=b_3\sqrt{\tfrac{1-C}{1+C}}E_3,\
\nabla_{E_3}E_2=-b_5\sqrt{\tfrac{1+C}{1-C}}E_3,\ \nabla_{E_3}E_3=-b_3\sqrt{\tfrac{1-C}{1+C}}E_1
+b_5\sqrt{\tfrac{1+C}{1-C}}E_2.
\end{aligned}\right.
\end{equation}

\begin{theorem}\label{thm:4.1abc}
Let $M$ be a hypersurface of $\mathbb{H}^2\times \mathbb{H}^2$ with constant sectional curvature $-\frac{1}{2}$.
If under the frame field $\{E_1,E_2,E_3\}$ defined by \eqref{eqn:2.5}, it holds
$$
C>0,\ b_1=\frac{1}{\sqrt{2}},\ b_4=-\frac{1}{\sqrt{2}},\ b_6=-\frac{C}{\sqrt{2}},\ b_2=b_3=b_5=0.
$$
%Then, $M$ is locally given by the immersion
%$\Phi:\Omega\subset\mathbb{R}^3\rightarrow\mathbb{H}^2\times\mathbb{H}^2$:
%$(t,r,s)\rightarrow(p(t,r),q(t,s))$, where $t<0$ and
%
%\begin{equation}\label{eqn:jr}
%\begin{aligned}
%&p(t,r)=\sqrt{\cosh(\sqrt{2}t)+1}(\cosh(r),\sinh(r),0)+\sqrt{\cosh(\sqrt{2}t)}(0,0,1),\\
%&q(t,s)=\sqrt{\cosh(\sqrt{2}t)}(1,0,0)+\sqrt{\cosh(\sqrt{2}t)-1}(0,\cos(s),\sin(s)),
%&p(t,r)=\Big(\sqrt{\cosh(\sqrt{2}t)+1}\cosh(r),\sqrt{\cosh(\sqrt{2}t)+1}\sinh(r),\sqrt{\cosh(\sqrt{2}t)}\Big),\\
%&q(t,s)=\Big(\sqrt{\cosh(\sqrt{2}t)},\sqrt{\cosh(\sqrt{2}t)-1}\cos(s),\sqrt{\cosh(\sqrt{2}t)-1}\sin(s)\Big),
%\end{aligned}	
%\end{equation}
Then, $M$ is an open part of Example \ref{exam:4.1ss1ss11ab1}.
\end{theorem}

\begin{proof}
By $b_3=b_5=0$, $b_6=-\frac{C}{\sqrt{2}}$ and $\nabla C=-2AX$, we have that product angle function satisfies
$$
E_1C=E_2C=0,\ E_3C=C\sqrt{2(1-C^2)}.
$$
Now, we choose two appropriate non-zero functions $\rho_1$ and $\rho_2$, such that
the new frame
$$
\{X_1=\rho_1E_1,\ X_2=\rho_2E_2,\ X_3=E_3\}
$$
satisfies that $[X_1,X_2]=[X_1,X_3]=[X_2,X_3]=0$. Specifically, by using
$[E_i,E_j]=\nabla_{E_i} {E_j}-\nabla_{E_j} {E_i}$, \eqref{eqn:LL} and $b_2=b_3=b_5=0$, we get
$$
[E_1,E_2]=0,\ [E_1,E_3]=-\sqrt{\tfrac{1-C}{2(1+C)}}E_1,\ [E_2,E_3]=-\sqrt{\tfrac{1+C}{2(1-C)}}E_2.
$$
Therefore
$$
\begin{aligned}
{\text [X_1,X_3]}&=-(\sqrt{\tfrac{1-C}{2(1+C)}}\rho_1+E_3\rho_1)E_1.
\end{aligned}
$$
Then $[X_1,X_3]=0$ is equivalent to
%
\begin{equation}\label{eqn:m1}
E_3\rho_1=-\sqrt{\tfrac{1-C}{2(1+C)}}\rho_1.
\end{equation}
Similarly, we get
$$
{\text [X_2,X_3]}=(-\sqrt{\tfrac{1+C}{2(1-C)}}\rho_2-E_3\rho_2)E_2,\ \
{\text [X_1,X_2]}=\rho_1(E_1\rho_2)E_2-\rho_2(E_2\rho_1)E_1.
$$
Then $[X_2,X_3]=[X_1,X_2]=0$ are equivalent to
%
\begin{equation}\label{eqn:m2}
E_3\rho_2=-\sqrt{\tfrac{1+C}{2(1-C)}}\rho_2,
\end{equation}
%
\begin{equation}\label{eqn:dm1dm4}
\rho_1 E_1\rho_2=\rho_2E_2\rho_1=0.
\end{equation}

Let $\rho_1=\sqrt{\frac{1+C}{2C}}$, $\rho_2=\sqrt{\frac{1-C}{2C}}$, then by $E_1C=E_2C=0$ and $E_3C=C\sqrt{2(1-C^2)}$,
we can check that \eqref{eqn:m1}--\eqref{eqn:dm1dm4} hold, which implies that $[X_1,X_2]=[X_1,X_3]=[X_2,X_3]=0$.

By the definition of the frame $\{X_1,X_2,X_3\}$ and using $[X_1,X_2]=[X_1,X_3]=[X_2,X_3]=0$,
we can identify $M$ with an open subset $\Omega_1$ of $\mathbb{R}^3$ locally and express the hypersurface $M$ by an immersion
$$
\begin{aligned}
\Phi_1:\Omega_1
\longrightarrow\mathbb{H}^2\times\mathbb{H}^2,\ \ \
(t_1,r,s)\mapsto (p(t_1,r,s),q(t_1,r,s)),
\end{aligned}
$$
such that $(\frac{\partial p}{\partial t_1},\frac{\partial q}{\partial t_1})=E_3$,
$(\frac{\partial p}{\partial r},\frac{\partial q}{\partial r})=\rho_1E_1$ and
$(\frac{\partial p}{\partial s},\frac{\partial q}{\partial s})=\rho_2E_2$.


By $C\neq0$, we integrate $E_3C=C\sqrt{2(1-C^2)}$ along the integral curves of $E_3$, and obtain
$$
C=\frac{1}{\cosh(\sqrt{2}t_1+\tilde{h})},
$$
where $\tilde{h}$ is a function on $M$ satisfying $E_3\tilde{h}=0$ and
$\sqrt{2}t_1+\tilde{h}<0$.


By $E_1C=E_2C=0$, $(\frac{\partial p}{\partial r},\frac{\partial q}{\partial r})=\rho_1E_1$ and
$(\frac{\partial p}{\partial s},\frac{\partial q}{\partial s})=\rho_2E_2$, we have that
$$
\frac{\partial C}{\partial r}=\frac{\partial C}{\partial s}=0\ \  {\rm and}\ \
\frac{\partial \tilde{h}}{\partial r}=\frac{\partial \tilde{h}}{\partial s}=0.
$$
It follows from $E_3\tilde{h}=0$ that $\tilde{h}$ is constant function. Thus, without loss of generality,
we can identify $M$ with an open subset $\Omega$ of $\mathbb{R}^3$ locally and express the hypersurface $M$ by an immersion
$$
\begin{aligned}
\Phi:\Omega
\longrightarrow\mathbb{H}^2\times\mathbb{H}^2,\ \ \
(t,r,s)\mapsto (p(t,r,s),q(t,r,s)),
\end{aligned}
$$
such that $(\frac{\partial p}{\partial t},\frac{\partial q}{\partial t})=E_3$,
$(\frac{\partial p}{\partial r},\frac{\partial q}{\partial r})=\rho_1E_1$,
$(\frac{\partial p}{\partial s},\frac{\partial q}{\partial s})=\rho_2E_2$ and
$C=\frac{1}{\cosh(\sqrt{2}t)}$, where $t=t_1+\frac{\tilde{h}}{\sqrt{2}}<0$. By $E_3=(\frac{\partial p}{\partial t},\frac{\partial q}{\partial t})$ and
$PN=CN+\sqrt{1-C^2}E_3$, we have $N=(\sqrt{\frac{1+C}{1-C}}\frac{\partial p}{\partial t},-\sqrt{\frac{1-C}{1+C}}\frac{\partial q}{\partial t})$.

From the definition of $P$  and using
$$
PE_1=E_1,\ \ PE_2=-E_2,
$$
we obtain that $dp,dq:T(\Omega)\rightarrow
T\mathbb{H}^2$ satisfy
\begin{equation}\label{eqn:5.49}
\left\{
\begin{aligned}
(dp(\tfrac{\partial}{\partial r}),0)&=\tfrac{1}{2}(d\Phi(\tfrac{\partial}{\partial r})+Pd\Phi(\tfrac{\partial}{\partial r}))=d\Phi(\tfrac{\partial}{\partial r}),\\
(0,dq(\tfrac{\partial}{\partial r}))&=\tfrac{1}{2}(d\Phi(\tfrac{\partial}{\partial r})-Pd\Phi(\tfrac{\partial}{\partial r}))=0,
\end{aligned}
\right.
\end{equation}
%
%\vskip-3mm
%
and
\begin{equation}\label{eqn:5.50}
\left\{
\begin{aligned}
(dp(\tfrac{\partial}{\partial s}),0)&=\tfrac{1}{2}(d\Phi(\tfrac{\partial}{\partial s})+Pd\Phi(\tfrac{\partial}{\partial s}))=0,\\
(0,dq(\tfrac{\partial}{\partial s}))&=\tfrac{1}{2}(d\Phi(\tfrac{\partial}{\partial s})-Pd\Phi(\tfrac{\partial}{\partial s}))=d\Phi(\tfrac{\partial}{\partial s}).
\end{aligned}
\right.
\end{equation}

The first equation of \eqref{eqn:5.50} shows that $p$ depends only on the $(t,r)$.
Similarly, from the second equation in \eqref{eqn:5.49} we derive that $q$ depends only on
the $(t,s)$. By $g(PE_3,E_3)=-C$, $AE_3=-\frac{C}{\sqrt{2}}E_3$ and $\nabla_{E_3}{E_3}=0$,
we have that
\begin{equation}\label{eqn:4.mm1}
\begin{aligned}
(\tfrac{\partial^2 p}{\partial t^2},\tfrac{\partial^2 q}{\partial t^2})&%=D_{E_3}{E_3}
=\bar{\nabla}_{E_3}{E_3}-g( D_{E_3}{E_3},\tfrac{(p,q)}{\sqrt{2}})\tfrac{(p,q)}{\sqrt{2}}
-g( D_{E_3}{E_3},\tfrac{(p,-q)}{\sqrt{2}})\tfrac{(p,-q)}{\sqrt{2}}\\
&=-\tfrac{C}{\sqrt{2}}N+\tfrac{1}{2}(p,q)-\tfrac{C}{2}(p,-q)\\
&=(-C\sqrt{\tfrac{1+C}{2(1-C)}}\tfrac{\partial p}{\partial t},C\sqrt{\tfrac{1-C}{2(1+C)}}\tfrac{\partial q}{\partial t})+(\tfrac{1-C}{2}p,\tfrac{1+C}{2}q),
\end{aligned}
\end{equation}
where $D$ is the canonical Euclidean connection on $\mathbb{R}_2^6$.
From $\frac{dC}{dt}=E_3C=C\sqrt{2(1-C^2)}$, we know that the solution of
\eqref{eqn:4.mm1} is given by
\begin{equation}\label{eqn:pq}
\begin{aligned}
p(t,r)=\sqrt{\tfrac{1+C}{C}}V_1(r)+\tfrac{1}{\sqrt{C}}V_2(r),\ \
q(t,s)=\tfrac{1}{\sqrt{C}}W_1(s)+\sqrt{\tfrac{1-C}{C}}W_2(s),
\end{aligned}
\end{equation}
where $V_1(r),V_2(r),W_1(s),W_2(s)\in \mathbb{R}_1^3$ are vector-value functions depending on variables $r$ and $s$. Thus, we have
$$
\begin{aligned}
E_1&=\tfrac{1}{\rho_1}(\tfrac{\partial p}{\partial r},\tfrac{\partial q}{\partial r})
%=\sqrt{\frac{2C}{1+C}}(\sqrt{\frac{1+C}{C}}\frac{d V_1(r)}{dr}+\frac{1}{\sqrt{C}}\frac{d V_2(r)}{dr},\vec{0})
=(\sqrt{2}\tfrac{d V_1(r)}{dr}+\sqrt{\tfrac{2}{1+C}}\tfrac{d V_2(r)}{dr},0),\\
E_2&=\tfrac{1}{\rho_2}(\tfrac{\partial p}{\partial s},\tfrac{\partial q}{\partial s})
%=\sqrt{\frac{2C}{1-C}}(\vec{0},\frac{1}{\sqrt{C}}\frac{d W_1(s)}{ds}+\sqrt{\frac{1-C}{C}}\frac{d W_2(s)}{ds})
=(0,\sqrt{\tfrac{2}{1-C}}\tfrac{d W_1(s)}{ds}+\sqrt{2}\tfrac{d W_2(s)}{ds}),\\
E_3&=(\tfrac{\partial p}{\partial t},\tfrac{\partial q}{\partial t})
=(-\sqrt{\tfrac{1-C}{2C}}V_1(r)-\sqrt{\tfrac{1-C^2}{2C}}V_2(r),
-\sqrt{\tfrac{1-C^2}{2C}}W_1(s)-\sqrt{\tfrac{1+C}{2C}}W_2(s)),
\end{aligned}
$$
and
$$
N=(\sqrt{\tfrac{1+C}{1-C}}\tfrac{\partial p}{\partial t},-\sqrt{\tfrac{1-C}{1+C}}\tfrac{\partial q}{\partial t})
=(-\sqrt{\tfrac{1+C}{2C}}V_1(r)-\tfrac{1+C}{\sqrt{2C}}V_2(r),
\tfrac{1-C}{\sqrt{2C}}W_1(s)+\sqrt{\tfrac{1-C}{2C}}W_2(s)).
$$



From the fact that $\langle p(t,r),p(t,r)\rangle=\langle q(t,s),q(t,s)\rangle=-1$, we have
$$
\begin{aligned}
&1+\langle V_1(r),V_1(r)\rangle+\tfrac{1}{C}(\langle V_1(r),V_1(r)\rangle+\langle V_2(r),V_2(r)\rangle)
+\tfrac{2\sqrt{1+C}}{C}\langle V_1(r),V_2(r)\rangle=0,\\
&1-\langle W_2(s),W_2(s)\rangle+\tfrac{1}{C}(\langle W_1(s),W_1(s)\rangle+\langle W_2(s),W_2(s)\rangle)
+\tfrac{2\sqrt{1-C}}{C}\langle W_1(s),W_2(s)\rangle=0.
\end{aligned}
$$
According to the independence of parameters $t,r$ and $s$, it follows from above two equations that
$$
\begin{aligned}
&\langle V_1(r),V_1(r)\rangle=-1,\ \langle V_1(r),V_2(r)\rangle=0,\ \langle V_2(r),V_2(r)\rangle=1.\\
&\langle W_1(s),W_1(s)\rangle=-1,\ \langle W_1(s),W_2(s)\rangle=0,\ \langle W_2(s),W_2(s)\rangle=1.
\end{aligned}
$$
Then, by $g(E_1,E_3)=g(E_2,E_3)=0$,
with the use of $\langle \frac{dV_1(r)}{dr},V_2(r)\rangle+\langle V_1(r),\frac{dV_2(r)}{dr}\rangle=0$
and $\langle \frac{dW_1(s)}{ds},W_2(s)\rangle+\langle W_1(s),\frac{dW_2(s)}{ds}\rangle=0$,
we can have
$$
\begin{aligned}
&\langle \frac{d V_1(r)}{dr},V_2(r)\rangle=\langle \frac{d W_1(s)}{ds},W_2(s)\rangle=0.
\end{aligned}
$$

Now, on one hand, for $V_1(r)\in \mathbb{H}^2$, if $V_1(r)$ is a constant vector, we have
$\frac{d V_1(r)}{dr}=0$ and $\frac{d V_2(r)}{dr}\neq0$. Then, it holds that
$$
\begin{aligned}
E_1&=\sqrt{\tfrac{2}{1+C}}(\tfrac{d V_2(r)}{d r},0),\\
D_{E_1}N&=\sqrt{\tfrac{2C}{1+C}}\tfrac{\partial N}{\partial r}=
\sqrt{\tfrac{2C}{1+C}}(-\tfrac{1+C}{\sqrt{2C}}\tfrac{d V_2(r)}{d r},0)=-\tfrac{1+C}{\sqrt{2}}E_1,
\end{aligned}
$$
which implies that $b_1=-\langle D_{E_1}N,E_1\rangle=\frac{1+C}{\sqrt{2}}$.
By the assumption that $b_1=\frac{1}{\sqrt{2}}$, we have $C=0$. It is a contradiction.
Thus, $V_1(r)$ is a curve of $\mathbb{H}^2$,
and $V_2(r)$ is a unit normal vector field of $V_1(r)\hookrightarrow\mathbb{H}^2$.
By taking re-parameterization to make $r$ being arc length parameter of $V_1(r)$, and assume that
$k(r)$ is the curvature of $V_1(r)$, we have
$$
\frac{d^2V_1(r)}{dr^2}=V_1(r)+k(r)V_2(r), \ \
\frac{dV_2(r)}{dr}=-k(r)\frac{dV_1(r)}{dr}.
$$
Now, we can have
$$
\begin{aligned}
E_1&=\Big((\sqrt{2}-k(r)\sqrt{\tfrac{2}{1+C}})\tfrac{d V_1(r)}{dr},0\Big),\\
D_{E_1}N&=\sqrt{\tfrac{2C}{1+C}}\tfrac{\partial N}{\partial r}=
((-1+k(r)\sqrt{1+C})\tfrac{d V_1(r)}{d r},0)=\tfrac{-1+k(r)\sqrt{1+C}}{\sqrt{2}-k(r)\sqrt{\tfrac{2}{1+C}}}E_1,
\end{aligned}
$$
which implies that $b_1=-\langle D_{E_1}N,E_1\rangle=\frac{1-k(r)\sqrt{1+C}}{\sqrt{2}-k(r)\sqrt{\frac{2}{1+C}}}$.
By the assumption that $b_1=\frac{1}{\sqrt{2}}$ and $C\neq0$, we have $k(r)=0$,
i.e., $V_1(r)$ is a geodesic of $\mathbb{H}^2$.

On the other hand, for $W_1(s)\in \mathbb{H}^2$, if $W_1(s)$ is not constant vector,
then $W_1(s)$ is a curve of $\mathbb{H}^2$,
and $W_2(s)$ is a unit normal vector field of $W_1(s)\hookrightarrow\mathbb{H}^2$.
By taking re-parameterization to make $s$ being arc length parameter of $W_1(s)$, and assume that
$\tilde{k}(s)$ is the curvature of $W_1(s)$, we have
$$
\frac{d^2W_1(s)}{ds^2}=W_1(s)+\tilde{k}(s)W_2(s), \ \
\frac{dW_2(s)}{ds}=-\tilde{k}(s)\frac{dW_1(s)}{ds}.
$$
So, it follows that
$$
\begin{aligned}
E_2&=\Big(0,(\sqrt{\tfrac{2}{1-C}}-\sqrt{2}\tilde{k}(s))\tfrac{d W_1(s)}{ds}\Big),\\
D_{E_2}N&=\sqrt{\tfrac{2C}{1-C}}\tfrac{\partial N}{\partial s}=
(0,(\sqrt{1-C}-\tilde{k}(s))\tfrac{d W_1(s)}{d s})=\tfrac{\sqrt{1-C}-\tilde{k}(s)}{\sqrt{\tfrac{2}{1-C}}-\sqrt{2}\tilde{k}(s)}E_2,
\end{aligned}
$$
which implies that $b_4=-\langle D_{E_2}N,E_2\rangle=\frac{-\sqrt{1-C}+\tilde{k}(s)}{\sqrt{\frac{2}{1-C}}-\sqrt{2}\tilde{k}(s)}$.
By the assumption that $b_4=-\frac{1}{\sqrt{2}}$, we have $\sqrt{1-C}-\sqrt{\frac{1}{1-C}}=0$.
It is a contradiction. Thus, $W_1(s)$ is a constant vector, we have
$\frac{d W_1(s)}{ds}=0$ and $\frac{d W_2(s)}{ds}\neq0$.
Then, we get
$$
\begin{aligned}
E_2&=(0,\sqrt{2}\tfrac{d W_2(s)}{d s}),\\
D_{E_2}N&=\sqrt{\tfrac{2C}{1-C}}\tfrac{\partial N}{\partial s}=
(0,\tfrac{d W_2(s)}{d s})=\tfrac{1}{\sqrt{2}}E_2,
\end{aligned}
$$
which implies that $b_4=-\langle D_{E_2}N,E_2\rangle=-\frac{1}{\sqrt{2}}$.

In conclusion, we know that $V_1(r)$ is a geodesic of $\mathbb{H}^2$ and $V_2(r)$ is a unit normal
vector field of $V_1(r)\hookrightarrow\mathbb{H}^2$;
$W_1(s)$ is a constant vector in $\mathbb{H}^2$ and $W_2(s)\in T^1_{W_1(s)}\mathbb{H}^2$.
By taking re-parameterizations of $r$ and $s$, and up to an isometry of $\mathbb{H}^2\times \mathbb{H}^2$, we can assume that
$$
\begin{aligned}
V_1(r)&=(\cosh(r),\sinh(r),0),\ \ V_2(r)=(0,0,1),\\
W_1(s)&=(1,0,0), \ \ W_2(s)=(0,\cos(s),\sin(s)).
\end{aligned}
$$
Thus, combining \eqref{eqn:pq}, $C=\frac{1}{\cosh(\sqrt{2}t)}$ and $t<0$, we exactly get the immersion \eqref{eqn:jr}.
%Conversely, it can be directly checked that hypersurface constructed by \eqref{eqn:jr}
%has constant sectional curvature.
\end{proof}

\end{document}
